\documentclass[a4paper,10pt]{article}
\usepackage[utf8]{inputenc}
\usepackage[T1]{fontenc}
\usepackage[serbian]{babel}
\usepackage{hyperref}
\usepackage{amsmath}
\usepackage{amssymb}
%opening
\title{DP procedura za iskaznu logiku}
\author{Nikola Kuburović, 1048/2024}

\begin{document}

\maketitle

\begin{abstract}
  DP (Davis-Putnam) procedura za iskaznu logiku ispituje zadovoljivost date formule iskazne logike.
  Sastoji se iz tri koraka:
  \begin{itemize}
    \item Unit propagation
    \item Pure literal
    \item Variable elimination
  \end{itemize}
  pri čemu je glavni korak upravo Variable elimination, dok prva dva služe da se poboljšaju performanse algoritma.
\end{abstract}

\section{Uvod}
Jedan od velikih problema u Automatskom rezonovanju jeste problem ispitivanja \emph{zadovoljivosti} formule:
za datu formulu, neophodno je ispitati da li postoji dodela istinitosnih vrednosti njenim promenljivim, takva da je formula tačna.

U našem slučaju, u pitanju je problem ispitivanja
zadovoljivosti formule iskazne logike.

Davis-Putnam (DP) procedura, kojom se ovaj rad bavi,
je jedna od najranijih procedura za ispitivanje zadovoljivosti logičke formule, i nju su 1960. godine
predložili Martin Dejvis i Hilari Putnam \cite{davis_putnam}

Interesantno je da je originalni DP algoritam zapravo namenjen za ispitivanje zadovoljivosti formule logike prvog reda, i oslanja se na Herbrandovu teoremu i ideju
da se problem u logici prvog reda svede na ispitivanje
niza iskaznih formula, pri čemu se za te iskazne formule koristi upravo DP procedura koju ovaj rad opisuje.

Ostatak ovog rada organizovan je na sledeći način. U poglavlju \ref{sec:osnove} navodimo osnovne definicije, pojmove i notaciju koja će biti korišćena u ostatku rada. U poglavlju \ref{sec:opis_metode} dajemo detaljni opis algoritma koji je predmet ovog rada. Poglavlje \ref{sec:implementacija} sadrži detalje implementacije metode. Najzad, poglavlje \ref{sec:zakljucak} sadrži zaključna razmatranja i osvrt na metodu opisanu u radu.

\section{Osnove}
\label{sec:osnove}

\paragraph{Formule iskazne logike.} Grade se od skupa iskaznih \emph{promenljivih} (atoma) $p, q, r, \dots$, logičkih konstanti $\top$ (tačno) i $\bot$ (netačno), i logičkih veznika $\neg, \wedge, \vee, \Rightarrow, \Leftrightarrow$.
\paragraph{Valuacija.} Preslikavanje $v$ koje svakoj promenljivoj dodeljuje istinitosnu vrednost iz $\{\top, \bot\}$; ono se na uobičajen način proširuje na sve formule. 
\paragraph{Zadovoljivost.} Za formulu $F$ kažemo da je \emph{zadovoljiva} ako postoji valuacija $v$ takva da je $v(F) = \top$; u suprotnom je \emph{nezadovoljiva}.
\paragraph{Tautologija.}Formula je \emph{tautologija} ako je tačna pri svakoj valuaciji. \emph{SAT problem} jeste problem odlučivanja da li je data formula zadovoljiva.

 \paragraph{Literal.} Promenljiva ($p$) ili njena negacija ($\neg p$).
\paragraph{Klauza.} Disjunkcija literala, npr.\ $(p \vee \neg q \vee r)$.
\paragraph{Konjunktivna normalna forma.} Formula je u \emph{konjunktivnoj normalnoj formi} (KNF) ako je konjunkcija klauza. DP procedura radi isključivo nad formulama u KNF.

\paragraph{Uslovi zaustavljanja.}\label{par:uslovi_zaustavljanja} Uvodimo dve granične vrednosti koje čine uslove zaustavljanja algoritma:
\begin{itemize}
  \item \textbf{Prazna klauza} $\{\}$ predstavlja $\bot$: disjunkcija bez članova je netačna. Pojava prazne klauze znači da je formula \textbf{nezadovoljiva}.
  \item \textbf{Prazan skup klauza} $\{\}$ predstavlja $\top$: konjunkcija bez članova je tačna. Prazan skup klauza znači da je formula \textbf{zadovoljiva}.
\end{itemize}

\paragraph{Rezolucija.} Ako klauza $C_1$ sadrži literal $x$, a klauza $C_2$ literal $\neg x$, njihova \emph{rezolventa} po \emph{pivotu} $x$ je klauza
\[
  R = (C_1 \setminus \{x\}) \cup (C_2 \setminus \{\neg x\}).
\]
Rezolventa je logička posledica konjunkcije $C_1 \wedge C_2$. Rezolucija je pravilo izvođenja na kome počiva glavni korak DP procedure.

\paragraph{Tseytin transformacija.} Proizvoljna formula ne mora biti u KNF, a naivno prevođenje može eksponencijalno povećati broj klauza. \emph{Tseytin transformacija} zaobilazi taj problem tako što svakoj podformuli dodeljuje novu promenljivu i dodaje klauze koje iskazuju ekvivalenciju te promenljive sa odgovarajućom podformulom. Rezultat je formula u KNF koja je \emph{ekvizadovoljiva} sa polaznom (zadovoljiva ako i samo ako je polazna zadovoljiva), uz linearno povećanje veličine. U našoj implementaciji Tseytin transformacija je ulazni korak koji formulu prevodi u KNF pre pokretanja DP procedure.



\section{Opis metode}
\label{sec:opis_metode}

DP procedura ponavlja 3 glavna koraka dok ne dođe do odluke o tome da li je formula zadovoljiva ili ne:
\subsection{Unit propagation}
\emph{Jedinična klauza} je klauza sa tačno jednim literalom $l$. Da bi cela formula bila tačna, taj literal mora biti tačan. Zato:
\begin{itemize}
  \item svaka klauza koja sadrži $l$ je zadovoljena i uklanja se iz formule;
  \item iz svake klauze koja sadrži $\neg l$ uklanja se literal $\neg l$ (jer ne doprinosi istinitosti klauze).
\end{itemize}

Uklanjanje $\neg l$ može neku klauzu svesti na jediničnu, pa ponavljamo isti postupak, ili na praznu klauzu, a tada je formula \textbf{nezadovoljiva}.
Takođe, moguće je i da uklanjanjem klauza dobijemo prazan skup klauza, i tada zaključujemo da je formula \textbf{zadovoljiva}.

\subsection{Pure literal}
Literal $l$ je \emph{čist} (eng. \emph{pure}) ako se u celoj formuli pojavljuje samo u jednom polaritetu tj.\ $\neg l$ se ne javlja ni u jednoj klauzi. Tada uvek možemo tom literalu dodeliti vrednost koja ga čini tačnim, ne narušavajući nijednu drugu klauzu. Sve klauze koje sadrže $l$ postaju zadovoljene i uklanjaju se. Ovaj korak može uzrokovati zadovoljivost (ispražnjenjem skupa klauza), ali nikada nezadovoljivost, jer samo briše cele klauze, a nikada ne skraćuje klauzu, pa ne može stvoriti praznu klauzu (slučaj koji dovodi do UNSAT zaključka).

\subsection{Variable elimination}

Ovo je glavni, najvažniji korak DP procedure. Bira se \emph{pivot} promenljiva $x$ i skup klauza $S$ se podeli na tri dela:
\begin{itemize}
  \item $P$ :: klauze koje sadrže $x$;
  \item $N$ :: klauze koje sadrže $\neg x$;
  \item $R$ :: ostale klauze (ne sadrže $x$).
\end{itemize}
Zatim se generišu sve rezolvente po pivotu $x$: za \emph{svaki} par $(C_p, C_n) \in P \times N$ računa se rezolventa. Novi skup klauza je
\[
  S' = R \cup \{\, \text{rezolventa}(C_p, C_n) \mid C_p \in P,\ C_n \in N \,\},
\]
uz izbacivanje rezolventi koje su \emph{tautologije} (sadrže i $y$ i $\neg y$ za neku promenljivu $y$, pa su trivijalno tačne). Skup $S'$ ne sadrži više promenljivu $x$ i \emph{ekvizadovoljiv} je sa $S$ (predstavlja $\exists x.\, S$). Uslovi zaustavljanja su oni opisani u paragrafu \hyperref[par:uslovi_zaustavljanja]{Uslovi zaustavljanja}:
\begin{itemize}
  \item ako neka rezolventa bude prazna klauza $\Rightarrow$ \textbf{Formula je nezadovoljiva};
  \item ako $S'$ bude prazan skup klauza $\Rightarrow$ \textbf{Formula je zadovoljiva}.
\end{itemize}

\paragraph{Cena koraka i izbor pivota.} Eliminacija promenljive sa $|P| = a$ i $|N| = b$ pojavljivanja \emph{uklanja} $a + b$ klauza, ali \emph{dodaje} do $a \cdot b$ rezolventi, po jednu za svaki par iz Dekartovog proizvoda $P \times N$. Kada je $a \cdot b > a + b$, broj klauza raste, i to se kroz uzastopne eliminacije može akumulirati eksponencijalno. Ovo je razlog za nastanak \emph{DPLL procedure} par godina kasnije. Kako je broj generisanih rezolventi $a \cdot b$, mogli bismo da koristimo heuristiku za izbor pivota koja \emph{minimizuje} $a \cdot b$. Pronalaženje optimalnog redosleda eliminacije je NP-težak problem.
\\ \\
\emph{Napomena}: Naša implementacija ne koristi nikakvu heuristiku za odabir pivota, već najprostije samo uzme prvi dostupan literal iz prve klauze u skupu klauza, i tako u svakoj iteraciji.
\section{Implementacija}
\label{sec:implementacija}




U ovom poglavlju treba dati detalje implementacije. Treba navesti u kom programskom jeziku je metoda implementirana, kako je k\^od organizovan, koje su klju\v cne funkcije/klase/metode i koja je njihova uloga.

Takodje je bitno dati upustvo za prevodjenje i pokretanje projekta, kao i navesti softver koji je potreban za prevodjenje i pokretanje programa (prevodilac, alati, dodatne bibiloteke i sl.).

\section{Zaklju\v cak}
\label{sec:zakljucak}

% U zaklju\v cku treba dati osvrt na ceo rad, kao i na metodu koja je u njemu prikazana. Ceo rad ne treba da ima vi\v se od 10 strana. Nakon zaklju\v cka, navesti spisak relevantne literature, ako je ista kori\v s\' cena u izradi seminarskog. Svi navedeni radovi iz spiska moraju biti citirani negde u tekstu na ovaj na\v cin \cite{sat_handbook}. Primer spiska literature je
% dat u nastavku (primer knjige \cite{sat_handbook}, rada na konferenciji \cite{fast} i rada u \v casopisu \cite{dpll_t}).


\begin{thebibliography}{100}
\bibitem{davis_putnam}  Davis, Martin, and Hilary Putnam. \emph{A computing procedure for quantification theory.} Journal of the ACM (JACM) 7.3 (1960): 201--215.

\end{thebibliography}

\end{document}



% TODO: Dodaj Maric, Bankovic slajdove u thebibliography sekciju
